% =============================================================================
%  ci2lab — "Scaffolding, Not Fine-Tuning"
%  Working draft. Compile:  pdflatex main && bibtex main && pdflatex main x2
%
%  ---------------------------------------------------------------------------
%  STATUS OF THIS DRAFT (read before editing)
%  ---------------------------------------------------------------------------
%  Sections 1-3 (Introduction, Related Work, Methods) are written and are
%  grounded in verified citations and in a direct source audit of the two
%  baseline harnesses.
%
%  Section 4 (Results) IS A SCAFFOLD. The benchmark suite has been implemented
%  but HAS NOT YET BEEN RUN. Every number in Section 4 is a placeholder marked
%  \tbd. NOTHING in this draft should be read as an empirical finding until the
%  experiment matrix in Section 3.5 has actually been executed.
%
%  Do not remove the \tbd macros by hand. They disappear only when real numbers
%  from benchmarks/results/*.jsonl replace them.
% =============================================================================

\documentclass[11pt,a4paper]{article}

\usepackage[utf8]{inputenc}
\usepackage[T1]{fontenc}
\usepackage[english]{babel}
\usepackage{geometry}
\geometry{margin=2.5cm}
\usepackage{amsmath,amssymb}
\usepackage{booktabs}
\usepackage{graphicx}
\usepackage{xcolor}
\usepackage{listings}
\usepackage{multirow}

% --- Bibliography and hyperlinking -----------------------------------------
% PACKAGE ORDER MATTERS. natbib MUST be loaded BEFORE hyperref: hyperref patches
% natbib's \cite commands to make them clickable, and it can only do so if natbib
% is already defined. Loading hyperref first silently produces dead citations.
% Do not reorder these three lines.
\usepackage[numbers,sort&compress]{natbib}

% backref=page: each bibliography entry lists the page(s) that cite it, as a
% clickable link back into the text. Set backref=false to disable for camera-ready.
\usepackage[
  colorlinks=true,
  linkcolor=blue!50!black,   % internal links (sections, tables, equations)
  citecolor=blue!45!black,   % \citep / \citet -> jumps to the bibliography
  urlcolor=blue!55!black,    % \url and \href
  linktoc=all,
  backref=page,
  pdfusetitle
]{hyperref}

% Makes the doi={...} field printed by plainnat into a clickable link.
% MUST be loaded AFTER hyperref.
\usepackage{doi}

% Phrasing for the back-references printed in the bibliography.
\renewcommand*{\backref}[1]{}
\renewcommand*{\backrefalt}[4]{%
  \ifcase #1 %
    \relax
  \or
    {\footnotesize (cited on p.~#2)}%
  \else
    {\footnotesize (cited on pp.~#2)}%
  \fi
}

% --- Placeholder macro: every unmeasured quantity MUST use this. -------------
\newcommand{\tbd}{\textcolor{red}{\textbf{[TBD]}}}
\newcommand{\note}[1]{\textcolor{red!70!black}{\small\textsf{[#1]}}}

\title{\textbf{Scaffolding, Not Fine-Tuning:\\
Deterministic Harness Mechanisms for Local Agentic Systems\\
with Unreliable Tool Calling}}

\author{%
  \note{author list} \\
  \note{affiliation} \\
  \texttt{\note{contact}}
}
\date{\today}

\begin{document}
\maketitle

% =============================================================================
\begin{abstract}
% TARGET: 300 words. Current draft is ~295.
Research on tool-using language agents overwhelmingly assumes a capable
underlying model that emits well-formed, provider-native tool calls. Production
harnesses inherit this assumption: an audit of two widely used open-source
systems, DeepAgents and OpenCode, confirms that both bind tools exclusively
through native function-calling interfaces, with no recovery path when a model
emits a tool call as prose, and that neither inspects the host machine to decide
what it can run. This is a poor fit for local-first deployment, where hardware
dictates the model and the affordable models are precisely those whose
tool-calling is least reliable.

We take the opposite position. We treat imperfect model output as the normal
case and ask how much agentic reliability can be recovered by a deterministic
layer built \emph{around} an unmodified, off-the-shelf open-weight model --- with
no fine-tuning, no adapters, and no distillation. We present \textsc{ci2lab}, a
local-first agentic system whose harness (i) recovers tool calls from free text
through a multi-strategy parser cascade, (ii) derives its context window from
measured VRAM using a KV-cache cost model and propagates that single budget to
serving, compaction, and orchestration, and (iii) verifies agent claims against
environment state using a non-LLM evidence gate.

We evaluate on a hand-authored, contamination-resistant task suite graded by
hidden exit-code oracles, under a design that isolates the harness from the
model: the same open model is run under a competing harness, our single- and
multi-agent modes are compared, and each scaffolding mechanism is ablated
individually across a model-size ladder. Beyond task success we report a
false-success rate --- runs that claim completion the oracle refutes.

\note{Results sentence pending: the experiment matrix has been implemented but
not yet executed. No empirical claim is made in this draft.}
\end{abstract}

\medskip
\noindent\textbf{Keywords:} local-first AI agents; agent harness; scaffolding;
tool calling; small language models; on-device inference; context management;
deterministic verification; ablation study; privacy-preserving AI

% =============================================================================
\section{Introduction}
\label{sec:intro}

Language-model agents act on the world by emitting tool calls: structured
requests to read a file, run a command, search a corpus. The dominant research
programme treats the \emph{model} as the seat of agentic competence. When agents
fail, the field's reflex is to improve the model --- by fine-tuning it for tool
use, by distilling a prompted agent into a smaller network, or by training
specialised expert modules. This reflex is visible across the on-device
literature: AgentFlux trains LoRA adapters to decouple tool selection from
argument generation \citep{kadekodi2025agentflux}; ReST-style approaches distil
a prompted ReAct agent into a far smaller fine-tuned model; CAMPHOR builds a
hierarchy of fine-tuned small-model experts \citep{fu2025camphor}. Each is a
strong result. Each also assumes that the path to a reliable local agent runs
through modifying the model's weights.

Production harnesses encode a related but distinct assumption: that the model
already speaks the protocol. We audited the source of two widely used
open-source agent harnesses. DeepAgents \citep{deepagents} delegates its control
loop to a framework that binds tools via a native function-calling interface;
its own documentation states that a tool-calling model is a hard requirement,
and we found no fallback path anywhere in the package. OpenCode \citep{opencode}
drives generation through a streaming SDK and likewise consumes only structured,
provider-native tool calls; its sole repair mechanism renames a mistaken tool or
routes the call to an \texttt{invalid} sentinel. Neither system profiles the host
machine's memory or accelerators, and neither derives its context window from
anything other than model metadata.

For a cloud deployment, these assumptions are benign: frontier models emit
clean tool calls, and the context window is whatever the provider advertises.
For a \emph{local-first} deployment they are close to fatal, and for a reason
that is structural rather than incidental. On local hardware the machine chooses
the model. A 48\,GB accelerator admits a quantised 32B model; a laptop admits a
7B one; and the smaller the model, the less reliable its tool calling
\citep{lin2025robust,rabinovich2025robustness}. The very models a local-first
system is forced to run are the ones these harnesses handle worst. The
constraint compounds: weights and KV cache compete for the same memory, so the
hardware also bounds the context window, and therefore how much of a task the
agent can hold in view at once.

This inverts the problem. If the hardware fixes the model, and the model cannot
be improved without training it, then reliability must be recovered somewhere
else --- in the deterministic machinery that surrounds the model. We call this
machinery the \emph{harness}, and this paper asks a single question:

\begin{quote}
\emph{How much agentic reliability can a deterministic harness recover from an
unmodified, off-the-shelf local model --- and which mechanisms are responsible?}
\end{quote}

The question is not rhetorical, and the honest answer is not obviously in our
favour. There is published evidence on both sides. Supporting us: SWE-agent
showed that holding the model fixed and changing only the agent--computer
interface moves performance substantially \citep{yang2024sweagent}; CodeAct
showed that the \emph{format} of the action space alone shifts success rates by
up to twenty points \citep{wang2024codeact}; and Kapoor et al. argue that the
community has systematically conflated scaffold engineering with model
capability, reaching mistaken conclusions about the source of accuracy gains
\citep{kapoor2025aiagents}. Cutting against us: METR found that post-training
enhancements delivered a far larger gain than their own scaffolding effort on an
already agency-tuned model \citep{metr2024elicitation}, and at least one
factorial study reports that adding scaffolding components can \emph{reduce}
performance through cross-component interference \citep{liu2026moreisnotbetter}.
A paper that ignored this evidence would be advocacy, not science. We therefore
frame our claim narrowly and test it adversarially: we expect scaffolding
headroom to be largest exactly where METR's is smallest --- on small, local,
\emph{non}-agency-tuned models.

\paragraph{Contributions.}
\begin{enumerate}
  \item \textbf{A harness architecture for unreliable tool calling.} A
    multi-strategy parser cascade recovers tool calls that models emit as XML,
    fenced code, bare JSON, or prose; in our model catalogue, 82 of 87 entries
    require it (Sec.~\ref{sec:system-parsing}).
  \item \textbf{Hardware-derived context budgeting.} We size the context window
    from measured VRAM using a KV-cache cost model
    \citep{kwon2023pagedattention,ainslie2023gqa}, and propagate that single
    number to the serving parameter, the compaction thresholds, and the
    multi-agent chunk planner (Sec.~\ref{sec:system-context}). To our knowledge
    no prior agent system derives its context policy from the host machine.
  \item \textbf{Deterministic, non-LLM verification.} Agent claims are checked
    against environment state by plain code, not by a judge model
    (Sec.~\ref{sec:system-grounding}). A judge reading the transcript would be
    evaluating the very signal that a false success has already corrupted, under
    self-preference biases that are unavoidable when judge and agent share
    weights \citep{panickssery2024llmevaluators} --- as they must in a
    single-model local deployment.
  \item \textbf{A harness-isolating experimental design}, in which the model is
    held fixed while the harness varies, each mechanism is ablated
    independently, and the whole matrix is repeated across a model-size ladder
    to test whether scaffolding value grows as the model shrinks
    (Sec.~\ref{sec:methods}).
\end{enumerate}

% =============================================================================
\section{Literature Review}
\label{sec:related}

We organise the relevant work into three strata and then locate the gap between
them.

\subsection{The runtime stratum: making local inference possible}

A large body of systems work has made local inference practical on consumer
hardware. Post-training quantisation compresses weights to three or four bits
with modest degradation \citep{frantar2023gptq,lin2024awq}; efficient serving
runtimes manage the KV cache as the binding memory constraint
\citep{kwon2023pagedattention}; and empirical studies map the resulting
performance envelope across desktop, laptop, and edge devices
\citep{huang2025edge,martinez2025latency,xu2024ondevice}. This stratum answers
\emph{can the model run here}. It stops at tokens per second, and does not ask
what the agent above it should then do.

One result from this stratum is directly load-bearing for us. ACBench evaluates
compression's effect on \emph{agentic} capability rather than perplexity, and
finds that 4-bit quantisation costs only 1--3\% on tool use and workflow
generation, but 10--15\% on real-world application accuracy
\citep{dong2025acbench}. A quantised local model is therefore a legitimate
substrate for an agent --- but a fragile one, and fragile in exactly the
long-horizon dimension a harness is positioned to defend.

\subsection{The model stratum: making small models agentic by training them}

The on-device agent literature converges on a single strategy: modify the model.
AgentFlux decomposes tool calling into selection and argument generation and
trains separate LoRA adapters for each \citep{kadekodi2025agentflux}. CAMPHOR
fine-tunes a hierarchy of small-model experts for planning, retrieval, and tool
interaction \citep{fu2025camphor}. ReST-style self-improvement distils a
prompted ReAct agent into a model two orders of magnitude smaller
\citep{aksitov2023rest}. Hammer trains on-device models with function masking to
resist misleading tool names \citep{lin2025robust}. A recent NVIDIA position
paper argues the endpoint of this
trajectory explicitly: small models are sufficient for the repetitive, narrow
invocations that dominate agentic systems \citep{belcak2025slm}.

We agree with the premise and dispute the method. Every one of these approaches
requires a training pipeline, and therefore a per-model, per-tool-set cost that
must be paid again whenever either changes. None asks what a purely training-free
intervention would buy.

\subsection{The agent stratum: control policy and context}

Above the model, work on agent design has produced the loop we build on
\citep{yao2023react}, self-correction mechanisms
\citep{shinn2023reflexion,madaan2023selfrefine} --- along with the important
caveat that \emph{intrinsic} self-correction, without external feedback, fails
and can even degrade accuracy \citep{huang2024llmscannot} --- and multi-agent
orchestration frameworks \citep{wu2024autogen,hong2024metagpt}. Closer to our
concerns, adaptive context management compresses agent state to fit tight
on-device budgets \citep{vijayvargiya2025adaptive}, prompt compression reduces
token footprint \citep{jiang2023llmlingua}, and OS-inspired designs page history
in and out of a virtual context \citep{packer2023memgpt}. Utility-guided
orchestration treats respond/retrieve/tool-call/verify/stop as an explicit policy
problem \citep{liu2026utility}.

Two properties unite this stratum. It is \emph{hardware-agnostic}: the context
budget is a configuration number or a value read from model metadata, never a
quantity derived from the machine. And it assumes the tool calls arrive intact:
the control policy decides \emph{which} tool to call, never what to do when the
model's call is malformed.

\subsection{Scaffolding as a first-class variable}

The claim that the harness --- not only the model --- determines agent
performance has recently acquired direct evidence. SWE-agent's central result is
that an agent--computer interface designed for a language model, with the model
held constant, substantially improves task resolution
\citep{yang2024sweagent}. CodeAct shows that merely changing the action format
from JSON or text to executable code raises success by up to 20\%
\citep{wang2024codeact}. The \textsc{Agentless} pipeline of
\citet{xia2025demystifying} makes the inverse point: a fixed pipeline can
outperform an autonomous agent, so the \emph{structure}, not the agency, carries
the result. At the meta level, Kapoor et al. show that
agent benchmarks routinely conflate scaffold engineering with model capability
\citep{kapoor2025aiagents}, and the Holistic Agent Leaderboard elevates
``scaffold'' to an explicit experimental dimension alongside model and benchmark
\citep{kapoor2025hal}.

The closest published work to our thesis is a recent preprint showing that a
single 8B model, placed in three scaffolded roles at inference time and with no
training whatsoever, roughly doubles tool-use completion and surpasses a system
four times its size \citep{mcclendon2026threeroles}. We differ in three ways.
That work assumes native tool calling and studies role decomposition; we target
models with \emph{no reliable tool-calling interface at all} and study output
recovery. It does not connect the harness to hardware; our context policy is
derived from it. And it reports task completion; we additionally measure
\emph{false success}.

\subsection{The alternative solution: constrain the decoder}

An obvious objection to our approach is that malformed tool calls are a solved
problem: constrain generation so that invalid output is unrepresentable. Guided
decoding indexes a finite-state machine over the vocabulary at near-zero
per-token cost \citep{willard2023outlines}; production engines enforce
context-free grammars efficiently \citep{dong2025xgrammar}; grammar-constrained
decoding without finetuning can beat both unconstrained and task-specific
fine-tuned models \citep{geng2023grammar}; and the machinery is available
on-device, not only through hosted APIs \citep{llamacpp}.

We take this objection seriously, and the literature does \emph{not} support the
convenient conclusion that constraints simply hurt. The best-known negative
result --- that format restrictions degrade reasoning \citep{tam2024speak} ---
is substantially undercut by evidence that the observed loss stems from subword
misalignment in the enforcement \emph{implementation} rather than from constraint
itself, with a properly aligned enforcer matching or exceeding unconstrained
generation \citep{beurer-kellner2024guiding}; a complementary line of work shows
that grammar-constrained decoding distorts the model's distribution and offers a
principled correction \citep{park2024grammar}.

Our position is therefore not that constrained decoding is wrong, but that it is
\emph{orthogonal and insufficient} --- and this follows from what the mechanism
is, not from any contested empirical claim. A grammar constrains the token
sequence of a \emph{single generation} to a language. That is exactly the
guarantee it offers, and exactly its limit. It cannot decide which tool the
situation calls for, recognise that the same call has now been issued three times,
stop a tool that keeps failing the same way, decide what to evict when the context
fills, or establish whether an agent that reports success has actually changed the
world. Those are properties of a \emph{trajectory} and of \emph{environment
state}; a per-token vocabulary mask is not the kind of thing that can express
them. Constrained decoding and a deterministic harness therefore operate on
different objects, and the presence of one does not discharge the need for the
other.

There is a further consideration, which we raise as a hypothesis to be tested
rather than as a settled result. Enforcing validity guarantees only that output
\emph{parses}, not that it is \emph{right}; the natural worry is that at small
scale it converts a visibly malformed call --- which a harness can detect and
repair --- into a well-formed call with wrong arguments, which nothing downstream
can detect at all. This concern is consistent with the distributional distortion
that \citet{park2024grammar} characterise formally, and preliminary evidence at
sub-3B scale points the same way \citep{ray2026constrainttax}, though that
evidence is a single-author preprint and we do not rest any claim on it. The
appropriate response is to measure it ourselves. Accordingly we treat constrained
decoding as a \textbf{baseline arm} (Sec.~\ref{sec:conditions}) and report schema
validity and task correctness as \emph{separate} quantities, so that a condition
which improves the former while degrading the latter is visible rather than
hidden behind a single success number.

\subsection{Verification, false success, and the choice of oracle}

Agents over-claim, and the mechanism is well established rather than anecdotal:
RLHF-trained models systematically prefer responses that match what the reader
expects to hear over responses that are true \citep{sharma2024sycophancy}. An
agent that has been trained to be agreeable and is then asked whether it finished
the task is being asked precisely the question its training biases it to answer
``yes'' to. We call the resulting failure \emph{false success}: the agent reports
completion that the environment refutes. Recent work characterises the phenomenon
directly and at scale \citep{advani2026falsesuccess}, and shows relatedly that
outcome-only grading can conceal procedural violations behind apparently
successful runs \citep{cao2026corruptsuccess}; both are preprints, and we cite
them as corroboration rather than as foundations. We claim no novelty for the
phenomenon. Our contribution is to treat its rate as a measured property of the
\emph{harness}, reported alongside task success rather than folded into it.

Our reason for measuring it with code rather than with a judge model is again
structural. A false success is, by construction, a case where the agent's prose
and the environment's state disagree; the agent's report is fluent, confident,
and wrong. A judge that reads the transcript is therefore evaluating exactly the
signal that has already been corrupted, and it does so under known and relevant
biases: LLM judges exhibit verbosity and self-enhancement effects
\citep{zheng2023judging}, and they favour their own generations in proportion to
their ability to recognise them \citep{panickssery2024llmevaluators} --- a
self-preference confound that is unavoidable in a single-model local deployment,
where judge and agent would necessarily be the same weights. An oracle that
inspects state rather than prose is immune to all of this by construction, which
is why our core suite is graded by exit code (Sec.~\ref{sec:oracle}) and our
grounding gate verifies claims in plain code (Sec.~\ref{sec:system-grounding}).
That design follows the attribution standard that a claim counts only if it is
verifiable against an independent source \citep{rashkin2023attribution}, and the
precedent for deterministic, out-of-model privilege control
\citep{shi2025progent,ruan2024toolemu}. Empirical support that judges
underperform on this specific task is reported by \citet{advani2026falsesuccess};
our argument does not depend on that measurement.

\subsection{The gap}

Three strata; two silences. \textbf{No published work connects the runtime
stratum to the agent stratum} --- hardware determines what can run, but no agent
system lets the machine determine its context policy. And \textbf{no published
work asks how far training-free scaffolding alone carries an off-the-shelf local
model} whose tool calling is unreliable, since the on-device agent literature
reaches for training by default. This paper occupies that gap.

Finally, the local-first premise itself is not merely a convenience. Cloud code
assistants have been shown to leak sensitive information present in the code they
ingest \citep{niu2023codexleaks}; the broader privacy threat model is surveyed in
\citet{miranda2025privacy}; and regulatory erasure obligations are technically
near-impossible to honour once data reaches a hosted model
\citep{zhang2024rtbf}. Running the agent locally does not make it more accurate.
It changes the trade: less model capability, in exchange for data control and
auditability.

% =============================================================================
\section{Methods}
\label{sec:methods}

\subsection{Case study: the \textsc{ci2lab} system}
\label{sec:system}

\textsc{ci2lab} is a local-first agentic command-line system, written in Python,
comprising four layers connected by explicit typed contracts.

\paragraph{Hardware profiling.} The system scans host RAM, VRAM, and accelerator
(via \texttt{psutil}, \texttt{nvidia-smi}, and Apple \texttt{sysctl}) and derives
an \emph{inference budget}: on a discrete GPU, total VRAM less a fixed reserve;
on CPU or unified memory, a conservative fraction of available RAM. The host is
classified into a coarse tier (edge / workstation / enterprise).

\paragraph{Model routing.} A catalogue of 87 open-weight models, each annotated
with memory requirements, native context length, and tool-calling mode, is
filtered to those that fit the measured budget and scored on a weighted
combination of quality, speed, hardware fit, and context. The essential fact for
this paper is the catalogue's composition: \textbf{82 of 87 models do not support
native tool calling}, and must be driven by prompting.

\paragraph{Transport.} Inference runs behind a pluggable backend interface. The
default targets Ollama's native chat endpoint rather than its
OpenAI-compatible shim, because only the former honours the context-length
parameter \citep{ollama}; an OpenAI-compatible backend covers vLLM, LM~Studio,
and llama.cpp. Swapping the model or the server is a configuration change.

\subsubsection{The parser cascade}
\label{sec:system-parsing}

The agent runs a ReAct loop \citep{yao2023react}. Where our loop departs from
convention is in how a tool call is obtained from the model's reply. Rather than
requiring a structured \texttt{tool\_calls} field, the harness applies an ordered
cascade of recovery strategies and accepts the first that yields a well-formed
call:

\begin{enumerate}\itemsep2pt
  \item native structured calls, when the model and backend support them;
  \item XML-style blocks (\texttt{<tool\_call>}, \texttt{<invoke>}), including
        normalisation of a vendor-specific markup dialect;
  \item fenced code blocks tagged with a tool name, carrying a JSON body;
  \item JSON objects, fenced or bare, located by a recursive scan;
  \item a bare tool name followed by a JSON object;
  \item generic fenced blocks, where a shell-tagged block is interpreted as a
        shell call.
\end{enumerate}

A separate detector recognises \emph{tool-shaped prose} --- text that is
evidently an attempted call but parses as none of the above. Such a turn does not
become the final answer; it triggers a corrective nudge. This is the mechanism
that lets the 82 non-native models in the catalogue act as agents at all.

\subsubsection{Loop robustness}
\label{sec:system-robustness}

Four further mechanisms, all task-agnostic and all deterministic, defend the
loop. \emph{Loop detection} hashes the tool calls issued each round and
intervenes when a round repeats. A \emph{retry governor} enforces two limits: an
identical call is not executed twice, and a tool that fails three times with the
same normalised error class --- even under varying arguments --- terminates the
attempt rather than thrashing. A \emph{read cache} collapses redundant retrievals
and is invalidated by any successful mutation. Finally, an \emph{anchor} is
re-injected each round: a synthetic message restating the original request
together with a progress digest built \emph{mechanically} from an evidence
ledger, with no model call --- a working-memory prosthesis for models that lose
the thread. Consistent with the finding that intrinsic self-correction fails
without external signal \citep{huang2024llmscannot}, every one of these
mechanisms is driven by observed tool and environment state, never by the model's
own introspection.

\subsubsection{Hardware-derived context budgeting}
\label{sec:system-context}

The context window is not configured; it is computed. Memory that is not occupied
by weights is available for the KV cache, and the KV cache grows linearly with
context length \citep{kwon2023pagedattention}. The harness therefore solves for
the affordable context:
\begin{equation}
  L_{\max} \;=\; \min\!\left(
    L_{\text{native}},\;
    \left\lfloor
      \frac{\alpha \,\bigl(B - W\bigr)}{c_{\mathrm{kv}}}
    \right\rfloor
  \right),
  \label{eq:ctx}
\end{equation}
where $B$ is the measured inference budget, $W$ the model's weight footprint,
$\alpha$ a safety fraction, $L_{\text{native}}$ the model's advertised window,
and $c_{\mathrm{kv}}$ the per-token KV cost. In principle
$c_{\mathrm{kv}} = 2 \cdot n_{\text{layers}} \cdot n_{\text{kv}} \cdot d_{\text{head}} \cdot s_{\text{dtype}}$;
the factor $n_{\text{kv}}$ (rather than the full head count) is what makes modern
grouped-query models affordable at long context \citep{ainslie2023gqa,shazeer2019mqa}.
Our implementation currently approximates $c_{\mathrm{kv}}$ by a single constant
scaled by weight size.
\note{Calibrate this constant against measured VRAM occupancy per model, or
replace it with the exact per-architecture formula; report the residual.}

The resulting $L_{\max}$ is a single source of truth: it is sent to the inference
server as the context parameter, it sets the thresholds of the compaction ladder,
and it determines whether a multi-agent review task is feasible at all. We cap
below the advertised window deliberately, since effective context is known to
fall well short of the nominal figure \citep{hsieh2024ruler,liu2024lost}.

Compaction is layered cheapest-first: superseded file reads are pruned; at 65\%
occupancy old tool results are stubbed; only at 80\% is an LLM summarisation
invoked; a mechanical trim is the final fallback.

\subsubsection{Deterministic grounding}
\label{sec:system-grounding}

In its multi-agent mode, \textsc{ci2lab} routes work through capability-scoped
roles whose tool allowlists are enforced outside the model, in the spirit of
deterministic privilege control \citep{shi2025progent}. Role selection itself is
rule-based, with no model call.

Reviewer roles do not emit free prose. They emit \emph{findings}, each typed by
its evidence: a verbatim quotation from the source document, an asserted absence,
or an external citation. Each finding is then verified \emph{in plain code}: the
quotation must occur in the document; the absent terms must genuinely be absent;
the cited URL must actually have been fetched during the run. Findings are
bucketed as verified, refuted, needs-check, or quarantined. No judge model is
involved --- a choice justified in Sec.~\ref{sec:related} by the poor performance
of LLM judges at precisely this task \citep{advani2026falsesuccess}, and modelled
on the attribution standard of \citet{rashkin2023attribution}.

\subsection{Task suite}
\label{sec:tasks}

We hand-author every task. Public agentic benchmarks may appear in pretraining
data, and contamination invalidates exactly the comparison we need to make
\citep{sainz2023contamination}; a small bespoke suite is contamination-resistant
by construction. The cost is statistical power, which we address in
Sec.~\ref{sec:stats} and concede in Sec.~\ref{sec:limitations}.

The suite spans four families, chosen so that each targets a distinct failure
mode (Table~\ref{tab:tasks}). Every task is self-contained, requires no network,
is bounded by a round limit and a wall-clock cap, and is graded by an oracle the
agent never sees.

\begin{table}[t]
\centering
\small
\begin{tabular}{@{}llp{5.4cm}p{4.4cm}@{}}
\toprule
\textbf{ID} & \textbf{Family} & \textbf{Failure mode tested} & \textbf{Oracle} \\
\midrule
CLI-01 & CLI  & Tool selection under a large file: search vs.\ read-all (context blow-up) & Planted token in final answer \\
CLI-02 & CLI  & Execute, read stderr, attribute the failing step & Names failing command + exit code \\
QA-01  & Q\&A & Locate a symbol in an unseen project vs.\ hallucinate a plausible one & Correct file/function/type; read-only \\
QA-02  & Q\&A & Multi-hop trace across three files, in order & Ordered three-hop chain \\
BUG-01 & Bug  & Localise and patch a root cause without breaking regressions & FAIL\_TO\_PASS \emph{and} PASS\_TO\_PASS \\
BUG-02 & Bug  & Two-file interacting regression; a one-line patch is insufficient & FAIL\_TO\_PASS \emph{and} PASS\_TO\_PASS \\
FEAT-01& Gen  & Implement to a hidden specification, including edge cases & Hidden spec tests \\
SEC-01 & Safety & Refuse or block a destructive command under policy & Tool blocked by permission engine \\
\bottomrule
\end{tabular}
\caption{The task suite. Difficulty is deliberately spread: CLI-01 and QA-01 are
warm-ups, QA-02/BUG-01/FEAT-01 discriminate, BUG-02 is the stretch task.}
\label{tab:tasks}
\end{table}

\subsection{The oracle}
\label{sec:oracle}

Grading is by exit code, never by a judge model. After the agent halts, hidden
test files are injected into the workspace --- into paths the agent never saw
during the run --- and the verifier executes. For repair tasks this is a
two-part contract: the previously failing test must now pass (FAIL\_TO\_PASS)
\emph{and} a set of regression tests must still pass (PASS\_TO\_PASS), so that
deleting an assertion or weakening the test cannot score as a fix. The verifier
additionally fails any run that modified a forbidden path, detected via version
control, which closes the remaining tampering route. This design follows the
execution-graded tradition of SWE-bench \citep{jimenez2024swebench} while
avoiding its contamination exposure.

\subsection{Experimental conditions}
\label{sec:conditions}

The design isolates the harness from the model. Every condition below runs the
\emph{same} task prompts, byte for byte, with no per-system prompt tuning.

\begin{description}\itemsep2pt
  \item[H1 --- System comparison (context).] \textsc{ci2lab} on a local open
    model versus hosted frontier agents as shipped. Reported for context only;
    it confounds model with harness and we do not rest any claim on it.
  \item[H2 --- Harness isolation (primary).] \textsc{ci2lab} versus a competing
    harness \emph{on the same self-hosted open model} $M$. Model held constant;
    the delta is attributable to orchestration. This is the experiment SWE-agent's
    ACI result \citep{yang2024sweagent} suggests should be decisive.
  \item[H3 --- Internal control.] Single-agent versus multi-agent
    \textsc{ci2lab} on $M$. No competitor, no confound.
  \item[A --- Mechanism ablation (the core contribution).] Each scaffolding
    mechanism is disabled \emph{individually}, leaving all others intact:
    parser cascade, retry governor, loop detection, anchors and progress digest,
    compaction ladder, and nudges. We additionally report the all-off condition,
    so that the mechanisms' joint effect can be compared against the sum of their
    individual effects: scaffolding components need not compose additively, and
    may interfere \citep{liu2026moreisnotbetter}. Reporting only individual
    ablations would hide that.
  \item[C --- Constrained-decoding baseline.] A condition in which tool calls
    are enforced by grammar-constrained decoding rather than recovered by the
    parser cascade \citep{willard2023outlines,llamacpp}. This is the honest
    alternative to our central mechanism and must be measured, not argued away.
    We report schema validity and task correctness as \emph{separate} quantities.
    The two are logically distinct --- a call can parse perfectly and still name
    the wrong file --- and a condition that improves one while degrading the other
    is invisible in a single success number.
  \item[L --- Model ladder.] The full matrix is repeated across a size ladder
    (7B / 14B / 32B of one model family). The interaction term is the headline:
    \emph{does the value of scaffolding grow as the model shrinks?} This is also
    where we test our reconciliation with METR's contrary finding
    \citep{metr2024elicitation}, which was obtained on a large, agency-tuned
    model.
\end{description}

\subsection{Metrics}
\label{sec:metrics}

\paragraph{Task success (primary).} Pass@1 from the oracle exit code. With $k$
samples per (task, condition) we also report the unbiased Pass@$k$ estimator of
\citet{chen2021codex}.

\paragraph{False-success rate (primary, and the one we consider most
diagnostic).} The fraction of runs in which the agent asserted completion but the
oracle refuted it. This is the failure mode that matters most for an agent
operating unattended: a stalled agent is merely useless, whereas an agent that
reports success it did not achieve actively misleads. It is measured against
environment state --- necessarily so, since the agent's own report is the
artifact under suspicion. The phenomenon and its prevalence are characterised in
\citep{advani2026falsesuccess,cao2026corruptsuccess}.

\paragraph{Efficiency.} Input and output tokens, and \emph{tokens per solved
task}, which is where a lean harness can win even when raw success trails. Cost
in USD is derived uniformly from measured tokens via a dated price table, and is
explicitly an \emph{imputed} hosted-equivalent figure for local runs, never an
invoice.

\paragraph{Latency.} End-to-end wall-clock, reported as median and p90, with a
discarded warm-up run per (agent, model) so that model-load time does not
contaminate the timed runs.

\paragraph{Failure attribution.} Every run is classified as solved, wrong,
stalled (round limit), or crashed. A harness that stalls fails differently from
one that confidently answers wrong, and the ablations are interpretable only if
the two are separated.

\paragraph{Safety.} Count of blocked or violating tool invocations under the
permission engine (SEC-01, and as a secondary axis throughout).

\subsection{Statistical treatment}
\label{sec:stats}

We draw $k=5$ samples per (task, condition) and report 95\% confidence intervals
by percentile bootstrap (2000 resamples, seeded for replayability). Point
estimates without intervals are not reported: with a suite this size, the
interval is the result \citep{miller2024errorbars}.

We fix temperature at 0 and pin model digests and harness commits. We do
\emph{not} claim this makes runs deterministic --- batch-dependent kernel
non-invariance means even greedy decoding is not reproducible in general
\citep{he2025nondeterminism} --- which is precisely why repeated sampling and
intervals, rather than a single greedy run, are load-bearing here.

\subsection{Threats to validity}
\label{sec:threats}

\paragraph{Model asymmetry.} Any comparison against hosted frontier agents (H1)
measures model capability at least as much as harness quality. This is why H1
carries no claim, and why H2/H3/A hold the model fixed. This confound is exactly
the one \citet{kapoor2025aiagents} identify as endemic.

\paragraph{Competitor handicap.} Running a competing harness on an open model it
was not designed for may understate its quality. We state this plainly rather
than exploit it: the H2 result should be read as \emph{``how these harnesses
behave on the models a local-first user can actually run,''} not as a claim about
their engineering in general.

\paragraph{Small $N$.} Eight tasks cannot support a general claim about agentic
capability. We scope every conclusion to this suite.

\paragraph{Single hardware point.} The context-budget mechanism is validated on
one machine. Its generality is untested.
\note{Consider a second hardware point (e.g. a 16\,GB consumer GPU) --- it would
strengthen the hardware-derived-context claim considerably.}

\paragraph{Pre-registration.} The protocol, the task suite, and the metric
definitions are frozen before the matrix is executed, and reported exactly as
specified here. No post-hoc task selection.

% =============================================================================
\section{Results}
\label{sec:results}

\begin{center}
\fbox{\parbox{0.92\linewidth}{\small\textcolor{red}{\textbf{DRAFT NOTICE.}
The benchmark harness described in Sec.~\ref{sec:methods} is implemented, but
\textbf{the experiment matrix has not yet been executed}. Every quantity in this
section is a placeholder (\tbd). The tables below fix the \emph{shape} of the
result and the analysis that will be applied to it; they contain no data. No
empirical claim in this paper is currently supported.}}}
\end{center}

\subsection{H2: harness isolation on a fixed model}

\begin{table}[h]
\centering
\small
\begin{tabular}{@{}lccccc@{}}
\toprule
\textbf{Harness} & \textbf{Pass@1 (95\% CI)} & \textbf{Pass@5} &
\textbf{False success} & \textbf{Tokens / solved} & \textbf{Latency p50 (s)} \\
\midrule
\textsc{ci2lab} (single) & \tbd & \tbd & \tbd & \tbd & \tbd \\
\textsc{ci2lab} (multi)  & \tbd & \tbd & \tbd & \tbd & \tbd \\
Competing harness        & \tbd & \tbd & \tbd & \tbd & \tbd \\
\bottomrule
\end{tabular}
\caption{H2/H3. Same open model $M$ under different harnesses. The comparison is
meaningful only because the model is held constant.}
\label{tab:h2}
\end{table}

\noindent\emph{Analysis to apply.} If the interval on the \textsc{ci2lab}
$-$ competitor difference excludes zero, the harness contributes independently of
the model, and the direction of the effect answers the paper's central question
for model $M$. If it does not, we say so.

\subsection{Ablation: which mechanism carries the weight?}

\begin{table}[h]
\centering
\small
\begin{tabular}{@{}lccc@{}}
\toprule
\textbf{Condition} & \textbf{Pass@1 (95\% CI)} & \textbf{$\Delta$ vs.\ full} &
\textbf{False success} \\
\midrule
Full harness              & \tbd & --- & \tbd \\
$-$ parser cascade        & \tbd & \tbd & \tbd \\
$-$ retry governor        & \tbd & \tbd & \tbd \\
$-$ loop detection        & \tbd & \tbd & \tbd \\
$-$ anchors / digest      & \tbd & \tbd & \tbd \\
$-$ compaction ladder     & \tbd & \tbd & \tbd \\
$-$ nudges                & \tbd & \tbd & \tbd \\
\midrule
All mechanisms off        & \tbd & \tbd & \tbd \\
Constrained decoding (C)  & \tbd & \tbd & \tbd \\
\bottomrule
\end{tabular}
\caption{Mechanism ablation on model $M$. The ``all off'' row tests
super-/sub-additivity: if the joint effect exceeds the sum of individual effects
the mechanisms are complementary; if it falls short, they interfere
\citep{liu2026moreisnotbetter}.}
\label{tab:ablation}
\end{table}

\noindent\emph{Analysis to apply.} We expect the parser cascade to dominate on
fenced models and to be inert on native ones --- if it is \emph{not} inert on a
native model, our account of what it does is wrong, and we will say so. The
constrained-decoding row is the one we are least able to predict. The specific
question we put to it is whether it raises schema validity while \emph{worsening}
false success: that pattern would indicate constraint is converting detectable
failures into undetectable ones, whereas a uniform improvement would indicate our
parser cascade should simply be replaced by a grammar where one is available.
Both outcomes are publishable; only one is comfortable.

\subsection{Model ladder: does scaffolding matter more as the model shrinks?}

\begin{table}[h]
\centering
\small
\begin{tabular}{@{}lccc@{}}
\toprule
\textbf{Model} & \textbf{Pass@1, full harness} & \textbf{Pass@1, all off} &
\textbf{Scaffolding gain} \\
\midrule
7B  & \tbd & \tbd & \tbd \\
14B & \tbd & \tbd & \tbd \\
32B & \tbd & \tbd & \tbd \\
\bottomrule
\end{tabular}
\caption{The headline interaction. A gain that \emph{increases} as parameters
decrease is the paper's central claim; a flat or inverted profile refutes it.}
\label{tab:ladder}
\end{table}

\noindent\emph{Analysis to apply.} This table is the paper's falsifiable core. A
monotonically increasing scaffolding gain as the model shrinks supports the
thesis and reconciles it with \citet{metr2024elicitation}, whose contrary result
was obtained on a large, agency-tuned model. A flat profile means scaffolding is
a constant additive benefit, which is a weaker but still publishable claim. A
gain that \emph{grows with} model size would refute our framing outright, and we
commit in advance to reporting it as such.

\subsection{Discussion}

\note{To be written against real data. The discussion must address, at minimum:
(1) whether the H2 delta survives its confidence interval;
(2) which mechanisms are load-bearing and which are decorative --- we should
expect some of our own machinery to turn out not to matter, and should say so;
(3) whether constrained decoding dominates the parser cascade, in which case the
honest conclusion is that our central mechanism should be replaced --- and we
will report that;
(4) the false-success axis, where we predict the largest harness effect and the
weakest competitor performance;
(5) reconciliation with METR's contrary elicitation finding via the ladder.}

% =============================================================================
\section{Limitations}
\label{sec:limitations}

Eight hand-authored tasks on one hardware configuration, with one model family,
cannot establish a general claim about agentic systems; every result is scoped to
this suite. The competing harness is evaluated outside its design envelope, which
likely understates it. Our KV-cache constant is an approximation pending
calibration. And the phenomenon of false success is not our discovery
\citep{advani2026falsesuccess,cao2026corruptsuccess}; our contribution is to
treat it as a property of the \emph{harness} and to show it can be reduced
without a judge model.

\paragraph{A note on evidentiary weight.} Parts of the literature we engage with
are recent preprints that have not been through peer review --- notably the
characterisations of false success \citep{advani2026falsesuccess,%
cao2026corruptsuccess}, the report of cross-component interference in scaffolding
\citep{liu2026moreisnotbetter}, and the evidence on validity--correctness
trade-offs under constraint \citep{ray2026constrainttax}. We have deliberately
structured our arguments so that none of them depends on those sources: each is
cited as corroboration for a position that is argued from mechanism or from
peer-reviewed work, and each of the corresponding empirical questions is one this
paper measures directly rather than inherits. Where a preprint's finding
conflicts with ours, our data should be preferred, and where it agrees, it should
be read as convergent rather than as support.

% =============================================================================
\section{Conclusion}
\label{sec:conclusion}

\note{To be written once results exist. The conclusion must not overstate: the
defensible claim is about small, local, non-agency-tuned models on this suite ---
not about agentic systems in general.}

% =============================================================================
\section*{Reproducibility}

The harness, the task suite with its hidden oracles, the adapters, the pinned
model digests, the price table, and the raw per-run results will be released.
Every figure in this paper is regenerable from \texttt{results.jsonl} by a single
command.

\bibliographystyle{plainnat}
\bibliography{references}

\end{document}
